% ladr-macros.tex: notation and environments for the Axler transcription.

% test builds of a single section set their counters only when this is true
\newif\iftestbuild
% Loaded after greenbook.sty.  Nothing in greenbook.sty is changed.

% ---------------------------------------------------------------- notation
\DeclareRobustCommand{\F}{\mathbb{F}}
\DeclareRobustCommand{\R}{\mathbb{R}}
\DeclareRobustCommand{\C}{\mathbb{C}}
\DeclareRobustCommand{\Q}{\mathbb{Q}}
\DeclareRobustCommand{\Z}{\mathbb{Z}}
\DeclareRobustCommand{\N}{\mathbb{N}}
\DeclareRobustCommand{\Lin}{\mathcal{L}}           % \Lin(V,W)
\DeclareRobustCommand{\Poly}{\mathcal{P}}          % \Poly_m(\F)
\DeclareRobustCommand{\Mat}{\mathcal{M}}           % \Mat(T)
\DeclareMathOperator{\nul}{null}         % \null is a TeX primitive
\DeclareMathOperator{\range}{range}
\DeclareMathOperator{\spn}{span}
\DeclareMathOperator{\adj}{adj}
\DeclareMathOperator{\Real}{Re}
\DeclareMathOperator{\Imag}{Im}
\DeclareMathOperator{\Rank}{rank}
\DeclareMathOperator{\dett}{det}
\DeclareMathOperator{\sgn}{sign}
\renewcommand{\rank}{\operatorname{rank}}
\renewcommand{\tr}{\operatorname{tr}}
\renewcommand{\Re}{\operatorname{Re}}
\renewcommand{\Im}{\operatorname{Im}}
\renewcommand{\dim}{\operatorname{dim}}
% physics' \det reads a (...) argument, which breaks inside theorem notes
\renewcommand{\det}{\operatorname{det}}
\renewcommand{\ip}[2]{\langle #1, #2 \rangle}
\renewcommand{\norm}[1]{\lVert #1 \rVert}
\renewcommand{\abs}[1]{\lvert #1 \rvert}
\newcommand{\ftimes}[2]{\underbrace{#1}_{#2}}  % x\cdots x with "m times" under it

% physics.sty turns \dd, \dv ... into commands; nothing of Axler needs them.
% Axler's l is a plain l: undo the \ell mapping of greenbook.sty.

% ------------------------------------------------- numbering as in Axler
% Every numbered item (definition, result, example, notation, figure) shares
% one counter per chapter: 1.1, 1.2, ...  Sections are 1A, 1B, ...
\renewcommand\thesection{\thechapter\Alph{section}}
% Axler's subsections are unnumbered (but stay in the contents)
\setcounter{secnumdepth}{1}
\makeatletter
% greenbook numbers definitions, theorems and examples per chapter, but each
% on its own counter: make definition and example use the theorem counter
\let\c@definition\c@theorem
\let\c@example\c@theorem
% Axler's results carry no name of their own: "Result" in the margin
\declaretheorem[style=hangingtheo,name=Result,sibling=theorem]{result}
\declaretheorem[style=hangingtheo,name=Result,numbered=no]{result*}
\declaretheorem[style=hangingtheo,name=Notation,sibling=theorem]{notation}
\declaretheorem[style=hangingtheo,name=Definition,numbered=no]{definition*}
\declaretheorem[style=thmremark,name=Remark,numbered=no]{remark*}
% a numbered figure takes its number from the same counter
\renewcommand\thefigure{\thetheorem}
% \axfigure: a numbered figure in Axler's counter.  Use as
%   \begin{figure}[htbp] ... \axcaption{text}\label{ax:6.9} \end{figure}
\newcommand\axcaption[1]{\refstepcounter{theorem}%
  \addtocounter{figure}{-1}\caption{#1}}

% The chapters use artbook's \marginnote[offset]{text}; the marginnote
% package (loaded by greenbook) takes \marginnote{text}[offset].  Accept both.
\let\gb@marginnote\marginnote
\RenewDocumentCommand\marginnote{O{0pt} +m o}{%
  \IfNoValueTF{#3}{\gb@marginnote{\footnotesize #2}[#1]}%
                  {\gb@marginnote{\footnotesize #2}[#3]}}
\makeatother

% ------------------------------------------------- lists inside the text
% (a), (b), ... as Axler sets them; bullets keep the default.
\newlist{parts}{enumerate}{1}
\setlist[parts]{label=(\alph*),leftmargin=*,widest=m,itemsep=1pt,topsep=3pt}
\setlist[itemize]{itemsep=1pt,topsep=3pt,leftmargin=1.2em}
\setlist[enumerate]{itemsep=1pt,topsep=3pt}

% the named property lists of Axler (commutativity, associativity, ...):
%   \begin{axprops} \item[commutativity] text ... \end{axprops}
\newlist{axprops}{description}{1}
\setlist[axprops]{style=nextline,leftmargin=1em,font=\normalfont\bfseries,
  itemsep=2pt,topsep=3pt,parsep=0pt}

% a short attribution line for the solutions
\newcommand\solnote[2]{\par\smallskip\noindent{\bfseries #1.}\ {\itshape #2}\par}

% a small unnumbered figure in the label column, as Axler's margin figures:
%   \marginfig{<tikz or \includegraphics>}{caption}
\newcommand\marginfig[3][0pt]{\marginnote[#1]{\vskip0pt\centering #2\par\smallskip
  {\itshape #3}\par}}

% the quotation that closes some exercise sets
\newcommand\axquote[2]{\par\bigskip\begingroup\raggedleft\small\itshape
  \parbox{.8\linewidth}{\raggedleft #1\par\smallskip\upshape\textemdash #2}\par\endgroup}

% the arrows of the R^2 pictures
\tikzset{vec/.style={-{Stealth[length=2.2mm]},line width=.6pt},
         axis/.style={-{Stealth[length=1.6mm]},line width=.3pt,black!60}}

% exercise numbers restart in every section (1A.1, 1A.2, ...)
\xsimsetup{exercise/within=section}

% solutions end with the black square, as in the manual; \qedhere works in a
% closing display
\AtBeginDocument{\xsimsetup{solution/pre-hook={\pushQED{\qed}},
  solution/post-hook={\popQED}}}

% the small italic remark Axler sets under some exercises
\newcommand\exnote[1]{\par\smallskip{\small\itshape\leftskip=1.5em\relax #1\par}}
\let\axexnote\exnote
\let\exercisenote\exnote

% the text column is narrow: allow a little extra stretch before overfull lines
\setlength\emergencystretch{1.5em}

% same space after the last exercise of every section
\AddToHook{cmd/section/before}{\par\addvspace{\bigskipamount}}
\allowdisplaybreaks[1]

% each group of solutions gets its own section-level contents entry
\renewcommand{\printsolutionsbysection}{%
  \def\lastchapter{}\def\lastsection{}%
  \ForEachUsedExerciseByType{%
    \edef\currentchapter{\ExercisePropertyGet{##1}{##2}{chapter-value}}%
    \edef\currentsection{\ExercisePropertyGet{##1}{##2}{section-value}}%
    \ifx\lastchapter\currentchapter\else
      \section*{Chapter~\ExercisePropertyGet{##1}{##2}{chapter}}%
    \fi
    \ifx\lastsection\currentsection\else
      \subsection*{\S\ExercisePropertyGet{##1}{##2}{section}}%
      \addcontentsline{toc}{section}{Solutions to
        \ExercisePropertyGet{##1}{##2}{section}}%
    \fi
    \XSIMprint{solution}{##1}{##2}%
    \let\lastchapter\currentchapter
    \let\lastsection\currentsection}}
